

\section{Introduction\label{sec:introduction}}

We survey two related literatures: {\it dynamic discrete choice\/} (DDC) in structural econometrics (SE) and ``inverse reinforcement learning'' (IRL) in machine learning (ML). Both are focused on inferring preferences of individuals from observations of their choices over time 
assuming they are {\it rational\/}  i.e. their
choices maximize an expected reward (``utility'') function based on beliefs of how
their actions affect current and uncertain future states, outcomes, and payoffs, allowing  
for unobserved states (``noise'') that result in decision rules that 
appear to be probabilistic from the standpoint of an outside observer who only partially observes agents'
states, \citet{rust:1987}. The literatures differ in the methods used to
solve the decision problem and infer agents' preferences, and have different rationales for inferring/learning them.

SE focuses on  testing and evaluating economic theories and making {\it counterfactual predictions\/}
of the effect of changes in the environment and
economic policies on economic behavior, outcomes, and welfare.\footnote{\footnotesize In economics ``policy'' refers to government interventions (tax rates, regulations, etc.), whereas
in RL/MDP it refers to a decision rule mapping states to actions.} 
Structural models differ from reduced-form models that summarize behavior without inferring
preferences. Reduced-form models require historical policy variation for counterfactuals,
but cannot forecast unprecedented policies. Structural models simulate how agents would
react to new policies.
For example, Denmark has long had one of the world's highest tax rates on new cars, but with little historical variation 
that could be used by reduced form approaches to predict how reductions in the tax would affect the Danish economy.
\citet{GIMRS:2022} structurally estimated an equilibrium model of ownership and trading of cars to recover consumer preferences
under the high tax {\it status quo\/} regime, and used it to compute counterfactual equilibria under a variety of alterative tax rates.
They showed that by reducing the tax on new cars and increasing the gas tax, Denmark can: 1) improve the welfare of most of its citizens,
2) increase total tax revenue from cars, and 3) reduce total $\mbox{CO}_2$ pollution.
Credible structural models are a faster, cheaper way of predicting the effects of policy changes by 
``crash testing'' computer simulations of new policies before actually implementing them.

IRL \citep{ng:2000} builds on {\it reinforcement learning\/} (RL) \citep{BartoSutton:1998}. Its goal is to develop intelligent agents that can perform complex robotic and intellectual tasks without
the need for explicit programming or an explicit reward function to direct behavior. The success of ML 
is due to {\it supervised learning\/} where algorithms are trained to predict outcomes using labeled data. 
However supervised learning is data-intensive and does not naturally handle sequential decision making 
where actions have long-term consequences. RL overcomes this by training agents to maximize a reward through trial-and-error,
via repeated dynamic interactions with the environment. 
When the reward function is known, RL can  approximate optimal behavior without labeled examples. Chess is an example with a known reward (1 for a win, -1 for a loss and 0 for a draw). RL has been able to learn superhuman performance via {\it self-play,\/} i.e. repeatedly playing the algorithm against itself resulting in performance that exceeds the best human and computerized chess players, see e.g. \citet{Silver:2018}.

However, many AI and robotics tasks lack clear reward function definitions. \citet{christiano:2017} consider whether  ``to use reinforcement learning to train a robot to clean a table or scramble an egg. It's not clear how to construct a suitable reward function, which will need to be a function of the robot's sensors.''  
IRL solves this by learning a reward function from limited data on human or expert demonstrations. Using the inferred reward function, 
this enables further training more rapidly and cheaply via RL simulations resulting in
 intelligent behavior consistent with human preferences in situations outside of those in a limited demonstration
data set.\footnote{\footnotesize
IRL is closely related to Reinforcement Learning from Human Feedback (RLHF),
which trains agents using RL where it is unclear how to
specify a reward function, instead learning from human preferences \citep{christiano:2017}.}
For example, Google Maps trip data \citep{Barnes:2024} revealed travelers' route preferences reflecting trade-offs between traffic, distance, hills, safety, and scenery. They massively scaled IRL to estimate a route preference function (360M parameters, 110M trips) enabling RL to generate recommendations achieving 16-24\% improvement in global route quality, and to the best of our knowledge, represents the largest published study or IRL algorithms in a real-world setting to date.

Section \ref{sec:mdp} reviews MDPs, the common ancestor of SE and RL. Section \ref{sec:ddc} covers DDC estimation and identification. Section \ref{sec:irl} summarizes IRL. Section \ref{sec:conclusion} discusses differences and open problems.


%\begin{table}[h!]
%\centering
%\caption{List of abbreviations and acronyms}
%\begin{tabular}{|l|l|}
%\hline
%\textbf{Abbreviation} & \textbf{Full Name} \\
%\hline
%AI & Artificial Intelligence \\
%AIRL & Adversarial Inverse Reinforcement Learning\\
%AGI & Artificial General Intelligence \\
%AL & Apprenticeship Learning \\
%BC & Behavioral Cloning \\
%BIRL & Bayesian Inverse Reinforcement Learning\\
%CCP & Conditional Choice Probability \\
%CCS & Conditional Choice Simulator \\
%DDC & Dynamic Discrete Choice \\
%DM  & Decision Maker (Agent) \\
%DNN & Deep Neural Network \\
%DP  & Dynamic Programming \\
%DSE & Dynamic Structural Estimation \\
%DU & Discounted Utility \\
%EU   & Expected Utility\\
%IOC & Inverse Optimal Control \\
%IL & Imitation Learning  \\
%IRL & Inverse Reinforcement Learning\\
%MCMC & Markov Chain Monte Carlo \\
%MDP & Markovian Decision Process \\
%ML & Machine Learning \\
%MPE & Markov perfect Equilibrium \\
%MPEC & Mathematical Programming with Equilibrium Constraints\\
%NFXP & Nested Fixed Point \\
%%PG & Policy Gradient \\
%PI  & Policy Iteration \\
%%PPO & Proximal Policy Optimization \\
%RFE & Reduced Form Estimation \\
%RL & Reinforcement Learning  \\
%RLHF & Reinforcement Learning from Human Feedback \\
%RTDP & Real Time Dynamic Programming \\
%SA  & Successive Approximations \\
%TD  & Temporal Difference\\
%%TRPO & Trust Region Polcy Optimization \\
%%XAI, XRL & Explainable AI, RL \\
%\hline
%\end{tabular}
%\label{tab:abbrevs}
%\end{table}
